\documentclass[10pt,a4paper]{amsart}
\usepackage[margin=19mm]{geometry}
\usepackage{amsmath,amssymb,mathtools}
\usepackage[scr=rsfs,cal=euler]{mathalfa}
\usepackage{xfrac}
\usepackage[unicode,hidelinks,bookmarks=false]{hyperref}
\newcommand{\ie}{i.e.\ }
\newcommand{\cf}{cf.\ }
\newcommand{\resp}{respectively\ }
\newcommand{\Subsection}{\subsection}
\providecommand{\nobreakdash}{\nobreak\hbox{-}}
\setlength{\emergencystretch}{2em}
\multlinegap0pt
\usepackage{bm}
\usepackage{bbm}
\usepackage{dsfont}
\usepackage{xspace}
\usepackage{cancel}
\usepackage{mathtools}
\usepackage{amsthm}
\makeatletter
\@ifundefined{theorem}{\newtheorem{theorem}{Theorem}}{}
\@ifundefined{proposition}{\newtheorem{proposition}[theorem]{Proposition}}{}
\makeatother
\title[Binary Representation in Multicomplex and Clifford Algebras]{Binary Representation in Multicomplex and Clifford Algebras}
\author{Derek Courchesne, S\'ebastien Tremblay}
\date{Source: arXiv:2509.23369v1 (CC BY 4.0)}
\begin{document}
\maketitle

\noindent\emph{Source and licence.} Binary Representation in Multicomplex and Clifford Algebras. Version arXiv:2509.23369v1 (CC BY 4.0), distributed under the Creative Commons Attribution 4.0 International licence, as recorded in the arXiv licence metadata for this exact version and shown on the abstract page https://arxiv.org/abs/2509.23369v1. Full rights notice: licenses/arxiv-2509.23369-CC-BY-4.0.txt.\par\smallskip

\noindent\textit{Editorial scope.} Counted unit: the Theorem whose source label is \texttt{theorem : algorithm-like diagonal basis} in the article (source0.tex lines 691--696, source label \texttt{theorem : algorithm-like diagonal basis}), delivered together with the article's own complete original proof of it (source0.tex lines 697--760, one \texttt{proof} environment of 64 lines). No mathematics is written, repaired or re-derived: statement and proof are the article's own text line for line. Required context retained uncounted: XOR algebra : binary representation (lines 366--369); XOR algebra : values of s (lines 402--410); hadamard-solution matrix def and properties (lines 674--678). Every definition, display and label of those lines is retained unchanged. Omitted from this package and not redistributed: 1--365, 370--401, 411--673, 679--690, 761--931. Nothing inside a retained excerpt is omitted or altered: every retained body is delivered exactly as the article's lines are written, so no omission receipt of any kind accompanies this package. Citations: the retained excerpts carry no citation command at all, so no bibliography entry of the article is transcribed and no reference is redistributed. Reference adaptation: none was needed and none was made; every reference inside the delivered unit resolves inside it, so no unresolved reference or citation is produced. Printed slips kept literal: no printed word, symbol, hypothesis or number of the counted statement or of its proof was altered. Layout and class: the article's own class is \texttt{article} with packages including amsfonts, amsmath, amssymb, amsthm; the wrapper is the lane's pinned \texttt{amsart} editorial wrapper at 10pt on A4, which restates the theorem-like environments the retained text needs and adds the article's own macro definitions that the retained text needs (none), with extra wrapper packages (bm, bbm, dsfont, xspace, cancel, mathtools). The delivered numbering is the unit's own. Assets: no retained span contains a figure environment, a graphics-inclusion command, a PDF-page inclusion, a table or a TikZ picture; no image file is copied into this package, the package declares no asset dependency and no image file is redistributed with it. Licence: version arXiv:2509.23369v1 of this article is distributed under the Creative Commons Attribution 4.0 International licence, as recorded in the arXiv licence metadata for this exact version and shown on the abstract page https://arxiv.org/abs/2509.23369v1. A full rights notice is delivered at licenses/arxiv-2509.23369-CC-BY-4.0.txt. Characters: every retained excerpt is pure ASCII, so the delivered engine, XeLaTeX with its default Latin Modern text font, reports no missing character.
\begin{equation}
    \label{XOR algebra : binary representation}
    \mathrm e_p := \prod_{k=0}^{n-1} u_{k+1}^{p_k}, \quad u_j^0 = 1, \quad j = 1,\ldots,n.
\end{equation}

\begin{proposition}
    \label{XOR algebra : values of s}
    For any hypercomplex algebra, the multiplier function for the standard basis has the following properties:
    \begin{enumerate}
        \item $s(p,q) \in \{ -1,0,+1 \}, \forall p,q$;
        \item if $u_i^2 \in \{ -1,+1 \}, \forall i$, then $s(p,q) \in \{ -1, +1 \}, \forall p,q$;
        \item if $u_i^2 =1, \forall i$ and $\lambda=1$, then $s(p,q) =1, \forall p,q$.
    \end{enumerate}
\end{proposition}

\begin{proposition}
    \label{hadamard-solution matrix def and properties}
    Let $T_n := [\nu_q (-1)^{\langle p,q\rangle}]_{p,q=0}^{2^n-1}$ where $\nu_q \in \mathbb{C}$ and $|\nu_q| = 1$. Then
    $$ T_n^{\dagger} T_n = 2^n I_{2^n}. $$
\end{proposition}

\begin{theorem}
    \label{theorem : algorithm-like diagonal basis}
    Let $s$ be the multiplier function of a commutative hypercomplex algebra $\mathcal{A}$ with $n \geq 0$ principal units $\{1,u_1, \ldots,u_n\}$ such that $u_i^2 \neq 0, \forall i$. Then there is a multiplicative diagonal basis given by the change-of-coordinates matrix
    $$ T_n := [\nu_q (-1)^{\langle p,q\rangle}]_{p,q=0}^{2^n-1}, $$
    from the standard basis to the new basis, with $\nu_q^2 = s(q,q)$.
\end{theorem}

\begin{proof}
    Let's denote by $\mathrm e_p, p=0,\ldots,2^n-1$ the elements of the standard basis as defined in equation (\ref{XOR algebra : binary representation}). From proposition \ref{XOR algebra : values of s}, $s(p,q) \in \{ -1, +1 \}$ since $u_k^2 \in \{ -1, +1 \}$ which implies that we can choose $\nu_q \in \{ 1,i \}$ with $\nu_q^2 = 1^2 = 1$ for the cases where $s(q,q) = 1$, and $\nu_q^2 = i^2 = -1$ for the cases where $s(q,q)=-1$. 
    
    % Note that since all the units of a multiperplex space square to $+1$, $s(q,q) = 1$ for all $q$ and always have a root in $\mathbb{R}$.
    
    We also denote by $\tilde{\mathrm{e}}_p, p=0,\ldots,2^n-1$ the elements of the new basis using the change-of-coordinates matrix $T_n$:
    $$ \tilde{\mathrm e}_k := \sum_{\alpha=0}^{2^n-1} (T_n)^{-1}_{\alpha,k} \mathrm e_\alpha. $$
    Since $|\nu_q| = 1$, we know from proposition \ref{hadamard-solution matrix def and properties} that the inverse of $T_n$ exists and is given by
    $$ (T_n)^{-1} = \frac{T_n^{\dagger}}{2^n}. $$
    The multiplier $\tilde{s}$ and the index function $\tilde{r}$ of the new basis are defined such that
    $$ \tilde{\mathrm e}_p \tilde{\mathrm e}_q = \tilde s(p,q) \tilde{\mathrm e}_{\tilde r(p,q)} $$
    and this basis is diagonal if $\tilde{s}(p,q) = \delta_{p,q}$ and $\tilde{r}(p,q) = p \delta_{p,q}$. We then proceed by induction, noting that the diagonal basis trivially exists when $n=0$. Suppose that the diagonal basis $\{ \hat{\mathrm e}_p \}_{p=0}^{2^n-1}$ exists for an algebra with $n$ principal units, obtained with the matrix $T_n$, and we add the unit $u_{n+1}$ such that $u_{n+1}^2 \in \{ -1, +1 \}$. Then for $p,q < 2^n$,
    \begin{equation}
        \label{multiperplex and multicomplex diagonal basis theorem : elements e}
        \mathrm e_{p+2^n} = \prod_{k=0}^{n} u_{k+1}^{(p+2^n)_k} = \prod_{k=0}^{n-1} u_{k+1}^{p_k} \cdot u_{n+1} = u_{n+1} \mathrm e_p.
    \end{equation}
    By definition of the bitwise scalar product,
    \begin{equation}
        \label{multiperplex and multicomplex diagonal basis theorem : bitwise scalar product}
        \begin{aligned}
            & \langle p + 2^n,q\rangle = \langle p,q + 2^n\rangle =  \langle p,q\rangle,\\
            & \langle p + 2^n,q + 2^n\rangle = \langle p,q\rangle + 1.
        \end{aligned}
    \end{equation}
    And by definition of $\nu_q$,
    $$ s(q+2^n, q+2^n) = \mathrm e_{q+2^n} \mathrm e_{q+2^n} = u_{n+1}^2 \mathrm e_q \mathrm e_q = s(2^n, 2^n) s(q,q) $$
    \begin{equation}
        \label{multiperplex and multicomplex diagonal basis theorem : elements nu}
        \Rightarrow \nu_{q+2^n} = \sqrt{s(q+2^n, q+2^n)} = \sqrt{s(2^n, 2^n) s(q,q)} = \nu_{2^n} \nu_q.
    \end{equation}
    From these last equations, we get
    $$
        \begin{aligned}
            (T_{n+1})_{p,q} &= \nu_q (-1)^{\langle p,q\rangle} = (T_n)_{p,q},\\
            (T_{n+1})_{p+2^n,q} &= \nu_q (-1)^{\langle p+2^n,q\rangle} = \nu_q (-1)^{\langle p,q\rangle} = (T_n)_{p,q},\\
            (T_{n+1})_{p,q+2^n} &= \nu_{q+2^n} (-1)^{\langle p,q\rangle} = \nu_{2^n}\nu_q (-1)^{\langle p,q\rangle} = \nu_{2^n}(T_n)_{p,q},\\
            (T_{n+1})_{p+2^n,q+2^n} &= \nu_{q+2^n} (-1)^{\langle p+2^n,q+2^n\rangle} = -\nu_{2^n}\nu_q (-1)^{\langle p,q\rangle} = -\nu_{2^n}(T_n)_{p,q},
        \end{aligned}
    $$
    or in the form of a block matrix
    \begin{equation}
        T_{n+1} = \begin{bmatrix}
            T_n & \nu_{2^n} T_n\\
            T_n & - \nu_{2^n} T_n
        \end{bmatrix} \Rightarrow (T_{n+1})^{-1} = \frac{1}{2} \begin{bmatrix}
            (T_n)^{-1} & (T_n)^{-1}\\
            \nu^*_{2^n} (T_n)^{-1} & - \nu^*_{2^n} (T_n)^{-1}
        \end{bmatrix}.
    \end{equation}
    The elements of the new basis are given, for $p<2^n$, by
    $$ \begin{aligned}
        \tilde{\mathrm e}_p &= \sum_{\alpha=0}^{2^{n+1}-1} (T_{n+1})^{-1}_{\alpha,p} \mathrm e_\alpha\\
        &= \sum_{\alpha=0}^{2^n-1} \left[ (T_{n+1})^{-1}_{\alpha,p} \mathrm e_\alpha+ (T_{n+1})^{-1}_{\alpha+2^n,p} \mathrm e_{\alpha+2^n} \right]\\
        &= \frac{1}{2} \sum_{\alpha=0}^{2^n-1} \left[ (T_{n})^{-1}_{\alpha,p} \mathrm e_\alpha+ \nu_{2^n}^* (T_{n})^{-1}_{\alpha,p} u_{n+1}\mathrm e_{\alpha} \right]\\
        &= \hat{\mathrm e}_p \frac{(1 + \nu^*_{2^n} u_{n+1})}{2}.
    \end{aligned} $$
    and in the same way,
    $$ \tilde{\mathrm e}_{p + 2^n} = \hat{\mathrm e}_p \frac{(1 - \nu^*_{2^n} u_{n+1})}{2}. $$
        From the following properties,
        $$ \left( \frac{1 + \nu^*_{2^n} u_{n+1}}{2} \right)^2 = \frac{(1 + \nu^*_{2^n} u_{n+1})}{2}, $$
        $$ \left( \frac{1 - \nu^*_{2^n} u_{n+1}}{2} \right)^2 = \frac{(1 - \nu^*_{2^n} u_{n+1})}{2}, $$
        $$ \frac{(1 + \nu^*_{2^n} u_{n+1})}{2} \cdot \frac{(1 - \nu^*_{2^n} u_{n+1})}{2} = 0, $$
        and the fact that $\{ \hat{\mathrm e}_p \}_{p=0}^{2^n-1}$ is a diagonal basis, it is clear that $\{ \tilde{\mathrm e}_k \}_{k=0}^{2^{n+1}-1}$ is a diagonal basis of the new algebra obtained by adding the unit $u_{n+1}$.
\end{proof}

\end{document}
